\documentclass[sigconf]{acmart}

\AtBeginDocument{%
  \providecommand\BibTeX{{Bib\TeX}}}

\title{Learned Index Structures: Recent Advances and Challenges}

\author{Barna Berekméri}
\affiliation{%
  \institution{Eötvös Loránd University (ELTE)}
  \city{Budapest}
  \country{Hungary}}
\email{barna.berekmeri@gmail.com}

\renewcommand{\shortauthors}{Barna Berekméri}

\begin{document}

\begin{abstract}
Index structures are fundamental to database performance. Recently, the
idea of replacing traditional data structures with machine learning
models has gained attention. This paper surveys recent progress in
learned indexes, focusing on learned Bloom filters, recursive models,
and multidimensional extensions.
\end{abstract}

\keywords{Learned index, Bloom filter, LSM-tree, multidimensional indexing, databases}

\maketitle

\section{Introduction}

Index structures are central to the performance of database management
systems, enabling efficient data access and query execution.
Traditional index structures, such as B-trees for range queries,
hash tables for point lookups, and Bloom filters for approximate
membership checks, have been refined over decades and serve as the
base of most commercial systems. These data structures provide
predictable worst-case guarantees and are well understood in terms of
their asymptotic complexity. However, the exponential growth of data
volumes, the increasing heterogeneity of workloads, and the demand for
low-latency data-intensive applications have exposed limitations of
classical index designs.

The fundamental challenge is that traditional indexes treat the key
distribution as unknown or adversarial, which means that they often
waste memory and compute resources to handle worst-case scenarios that
rarely occur in practice. In contrast, real-world data distributions
are typically skewed, correlated, or otherwise structured. This
observation has motivated the exploration of machine learning (ML) as
a tool to exploit such structure for indexing. The idea of
\emph{learned indexes} is that indexes can be replaced, or at least
augmented, by predictive models that approximate the mapping between
keys and their storage positions. If the model can be trained to
predict positions or membership efficiently, then the cost of lookups
and the memory overhead of index structures can be significantly
reduced \cite{kraska2018case}.

This idea represents a paradigm shift: rather than viewing an index as
a rigid data structure, we view it as a learned function that
leverages the statistical properties of the data. Over the past five
years, this line of research has generated intense interest in both
theory and system communities. Numerous papers in SIGMOD, VLDB, and
ICDE have explored new learned index structures, extensions to
multidimensional data, hybrid architectures that combine models with
traditional indexes, and practical evaluations in storage engines
\cite{howgood2024,fidalgo2025learned,sato2024fast,sato2025cascaded,canlearned2024}.

Despite promise, learned indexes also pose serious challenges.
They often lack worst-case guarantees, struggle with dynamic updates,
and can be costly to train or adapt when the data distribution
changes. Moreover, integrating learned indexes into storage engines
such as LSM trees or into disk-based systems raises new questions about
I/O complexity, cache behavior, and model management. In the
following, we provide background and review recent contributions in
this area.

\section{Background and Related Work}

\subsection{The Learned Index Paradigm}

The concept of learned indexes was introduced by Kraska et
al.~\cite{kraska2018case}, who observed that a key-to-position mapping
in a sorted array can be approximated by a cumulative distribution
function (CDF). In this view, a B-tree traversal can be replaced by
computing a model prediction and probing near the predicted location.
This led to recursive model indexes (RMIs), where the upper levels of the
hierarchy choose which model to use in the lower levels, mimicking the
branching behavior of a tree, but in functional form.

The original proposal demonstrated significant performance
improvements in lookup latency and memory footprint, sometimes by
orders of magnitude. However, subsequent studies also revealed
weaknesses: RMIs can suffer from high prediction errors in regions of
the data distribution, leading to expensive verification steps, and
their update performance is often inferior to that of B trees.

\subsection{Multidimensional Learned Indexes}

Extending learned indexes to multidimensional data is an active area
of research. Workloads such as spatial queries, multiattribute range
queries, or nearest-neighbor search demand indexes that operate over
two or more dimensions. Classical solutions include kd-trees, R-trees,
and their many variants. Learned approaches such as Flood, LISA, and
RSMI attempt to partition the multidimensional space and use models to
predict data placement.

The empirical study ``How Good Are Multidimensional Learned
Indexes?'' provides the most comprehensive evaluation to date
\cite{howgood2024}. The authors benchmark a wide range of learned and
traditional multidimensional indexes on synthetic and real-world data.
They find that learned indexes can outperform traditional structures
in terms of memory efficiency and range query latency, sometimes by up
to three orders of magnitude. However, the benefits are less clear for
updates and high-dimensional queries such as neighbors k nearest,
where traditional structures remain more robust. These findings
highlight both the potential and the current limitations of the learned
approaches.

\subsection{Learned Bloom Filters}

Another important line of research concerns Bloom filters, which are
probabilistic data structures for approximate membership queries.
Bloom filters are widely used in storage systems, including as part of
LSM-trees, to avoid unnecessary disk accesses. Traditional Bloom
filters guarantee no false negatives, but allow false positives with a
probability determined by their size and hash function choice.

The concept of learned Bloom filters replaces or augments the hash
functions with a classifier that predicts whether a key is likely to
exist in the set. A Bloom backup filter ensures the correctness by
capturing false negatives of the classifier
\cite{kraska2018case}. Recent advances include partitioned learned
Bloom filters (PLBFs), which allocate memory adaptively across
partitions of the key space. However, the original construction of
PLBFs was computationally expensive. Sato and Matsui propose faster
construction algorithms with theoretical guarantees
\cite{sato2024fast}. More recently, cascaded learned Bloom filters
(CLBFs) use multiple classifier stages to achieve better balance
between filter size and prediction accuracy, reducing memory usage by
about 24\% and accelerating negative lookups by up to 14$\times$
\cite{sato2025cascaded}.

The integration of learned Bloom filters into storage engines is
particularly noteworthy. Fidalgo and Ye propose two methods to embed
the Bloom filters learned into LSM trees \cite{fidalgo2025learned}. The
first method uses classifiers to skip certain filter checks, thereby
reducing latency, while the second replaces Bloom filters with
learned models plus small backup filters, saving 70--80\% of memory.
These experiments demonstrate substantial performance improvements
while preserving correctness guarantees.

\subsection{Construction Complexity and Maintenance}

While most early work on learned indexes focused on query performance,
a critical practical issue is the cost of constructing and maintaining
the models. Training a learned index on a large dataset can be
prohibitively expensive if not carefully optimized. The recent paper
``Can Learned Indexes be Built Efficiently?'' investigates this
problem in depth \cite{canlearned2024}. The authors analyze the
complexity of different model construction strategies and show how
sampling and incremental learning can mitigate training costs. They
argue that unless construction complexity is carefully managed, the
practical utility of learned indexes will be limited.

Another challenge is adapting to evolving data distributions. Static
models may degrade in accuracy as new data arrives or distributions
shift, leading to increased error rates and higher query costs. This
underscores the need for incremental training methods or hybrid
approaches that combine models with robust fallback mechanisms.

\section{Conclusion}

Learned indexes offer an exciting new perspective on a decades-old
problem in database systems. By leveraging machine learning to
approximate index functionality, they achieve significant gains in
memory and query performance. Recent work has extended the paradigm to
multidimensional queries and Bloom filters, and has begun to address
the practical challenges of construction complexity and integration
into storage engines. At the same time, open challenges remain in
ensuring robustness under updates, managing distribution drift, and
providing worst-case guarantees. These open problems define a rich
agenda for future research at the intersection of databases and
machine learning.


\bibliographystyle{ACM-Reference-Format}
\bibliography{references}

\end{document}
